\section{Discussion}

\noindent\textbf{Vulnerabilities in CNs}. Our research brings to light the pervasiveness and severity of iTrue vulnerabilities, which we attribute to a systemic specification-to-implementation gap. Requirements for protocol-syntax validation, protocol-semantics validation, and resource-availability checks are dispersed across 3GPP specifications, increasing the likelihood of incomplete enforcement in implementations and leading to recurring weaknesses such as iTrue. Crucially, these gaps threaten not only availability (e.g., denial of service) but also integrity (e.g., session hijacking), as demonstrated in our research, indicating that iTrue flaws are not isolated implementation bugs but systematic concerns.

Mitigating iTrue risks requires coordinated action across specifications, implementation, and deployment. At the specification level, 3GPP should update its threat model to reflect cloud-native deployments in which internal interfaces may be exposed, and should mandate mutual authentication on internal interfaces. At the implementation level, network-function developers should enforce strict syntactic and semantic validation at protocol boundaries, maintain uniqueness invariants for session and rule identifiers, and adopt bounded resource allocation with graceful rejection rather than process-ending assertions. At the deployment level, operators should adopt a zero-trust posture by removing unintended reachability, enforcing net-segmentation and least-privilege access between components, and continuously auditing cloud configurations to detect misconfigurations before they become exploitable.

We hope that the findings of this study will help catalyze such efforts, contributing to defense-in-depth protections for future CNs in particular, and carrier networks more broadly.

\vspace{2pt}\noindent\textbf{Limitations of iFinder}. 
\textit{iFinder} does not guarantee complete coverage of all iTrue vulnerabilities. The Discovery Agent is bounded by our seed pattern set, and vulnerabilities that do not match these patterns may be missed. Moreover, although our pipeline leverages procedure context and message schemas to reduce false positives, the LLM may still make incorrect feasibility judgments when exploitability hinges on subtle specification semantics, implementation-specific state transitions, or implicit invariants that are not captured by message structure alone. Future work should explore additional implicit assumptions underlying CNs and incorporate deeper, system-level analysis to further reduce false positives and broaden coverage.

\vspace{2pt}\noindent\textbf{Agent-based security analysis on carrier networks}. 
Our results demonstrate that multi-agent systems can automate vulnerability discovery across multiple protocols, achieving a level of scalability and accuracy that is difficult to attain with manual auditing or conventional, less intelligent tools. Of particular importance is the agentic system’s ability to perform multimodal analysis, jointly reasoning over code, specifications, and other artifacts, and to do so smoothly across different programming languages and protocols.

These capabilities are especially valuable for carrier-network security analysis, where the ecosystem is complex and heterogeneous, and where both the codebases and the supporting documentation are massive. Looking ahead, we expect agentic techniques to see broader adoption and tighter integration into pipelines for both specification development and system implementation, ultimately improving the security posture of this critical information infrastructure.
